% TOP_PROBLEM: {"id": 3700015, "title": "Zeros and one-points on three rays", "category_id": 8, "category": {"id": 8, "name": "analysis", "display_name": "Analysis", "description": "Limits, continuity, calculus, and function theory.", "slug": "analysis", "order_index": 8, "created_at": "2026-07-31T15:26:25.671Z"}, "status": "open", "problem_number": "AMR-036-0015", "set_id": 13, "statusQualification": "Restores the transcendental, infinitely-many-value-points setting implicit in the primary context; without it simple polynomial or omitted-value examples trivialize the catalogue wording."}
% FIRST_POSTED: 2026-10-01
% ENTRY_KIND: statement_only
% UnsolvedMath ID 3700015; source catalogue code AMR-036-0015
% !TeX program = lualatex
% Definition and sources only; this file is not an LLM solution attempt.
\documentclass[11pt,a4paper]{article}
\usepackage[margin=25mm]{geometry}
\usepackage{fontspec}
\setmainfont{lmroman10-regular.otf}[BoldFont=lmroman10-bold.otf,
 ItalicFont=lmroman10-italic.otf,BoldItalicFont=lmroman10-bolditalic.otf]
\usepackage{amsmath,amssymb,mathtools}
\usepackage{unicode-math}
\setmathfont{latinmodern-math.otf}
\AtBeginDocument{\let\setminus\smallsetminus}
\usepackage{xurl,hyperref}
\hypersetup{colorlinks=true,urlcolor=blue}
\setcounter{secnumdepth}{0}
\setlength{\parindent}{0pt}
\setlength{\parskip}{5pt}
\setlength{\emergencystretch}{3em}
\begin{document}
\section*{Zeros and one-points on three rays}\label{3700015}
{\small UnsolvedMath ID 3700015 · AMR-036-0015 · Analysis · AMR Open Problem Lists}
\subsection{Definitions and mathematical statement}
For a real angle $\theta$, write
$R_\theta=\{re^{i\theta}:r>0\}$ for the open ray from
the origin. A $1$-point of an entire function $f$ is a
solution of $f(z)=1$, counted with multiplicity when needed.
The intended problem concerns transcendental entire functions,
with infinitely many zeros and infinitely many $1$-points.

Does there exist such an $f$ and an angle
\[
       \frac{\pi}{3}<\alpha<\frac{\pi}{2},
                   \qquad\alpha\ne\frac{2\pi}{5},
\]
for which
\[
 \{z:f(z)=0\}\subset R_0,\qquad
 \{z:f(z)=1\}\subset R_\alpha\cup R_{-\alpha}?
\]
These are requirements on every zero and every $1$-point
in the entire plane, not only on their limiting directions
at large modulus. The two value sets may contain multiple
points; no simplicity condition is imposed. Since the
origin lies on none of the open rays, it is neither a
zero nor a $1$-point.

The target is existence at any angle in the displayed
range other than the distinguished angle, or a proof of
nonexistence throughout that range. No finite-order growth
assumption is imposed. Infinitely many preimages of each
value make the radial distribution condition substantive:
polynomials or functions omitting one of the two values
would permit vacuous or elementary configurations unrelated
to the intended three-ray problem.
\subsection{Short English statement}
Can infinitely many zeros and one-points of a transcendental entire function be confined to these three rays at another allowed angle?
\subsection{Sources}
\begin{itemize}
\item[S1] ULAM.AI and the UnsolvedMath Contributors, supplied problems.json, record 3700015 (AMR-036-0015), AMR Open Problem Lists. Catalogue identity and source transcription; CC BY 4.0.
\url{https://huggingface.co/datasets/ulamai/UnsolvedMath}.
\item[S2] Eremenko - My favorite unsolved problems, Radially distributed values, first question. Original source indexed by AMR-036-0015.
\url{https://www.math.purdue.edu/~eremenko/uns1.html}.
\item[S3] A. Eremenko, Radially distributed values. Author-maintained primary problem note.
\url{https://www.math.purdue.edu/~eremenko/dvi/radial.pdf}.
\end{itemize}
\end{document}
